\documentclass[12pt,oneside]{article}
\usepackage[utf-8]{inputenc}
\usepackage[margin=1in]{geometry}
\usepackage{setspace}
\usepackage{graphicx}
\usepackage{float}
\usepackage{booktabs}
\usepackage{amsmath}
\usepackage{amssymb}
\usepackage{natbib}
\usepackage{hyperref}
\usepackage{color}
\usepackage{soul}

% Line numbering for submission
\usepackage{lineno}
\linenumbers

% Double spacing
\doublespacing

% Define colors for notes
\definecolor{yellow}{RGB}{255,255,0}
\definecolor{cyan}{RGB}{0,255,255}

% Custom commands
\newcommand{\TODO}[1]{\textcolor{red}{[TODO: #1]}}
\newcommand{\CITE}{\textcolor{blue}{[CITE]}}
\newcommand{\FIG}[1]{Figure~\ref{fig:#1}}
\newcommand{\TAB}[1]{Table~\ref{tab:#1}}

%-----------
% TITLE PAGE
%-----------
\title{SPP1+ Tumor-Associated Macrophages and T Cell Exhaustion Predict Immunotherapy Response in Gastric Cancer: A Single-Cell Meta-Analysis}

\author{
    \TODO{Your Name}$^{1,2}$\thanks{Corresponding author: \textcolor{blue}{[email]}} \\
    \TODO{Co-Author 2}$^{1}$, \\
    \TODO{Co-Author 3}$^{1}$
}

\date{}

%-----------
% AFFILIATIONS (footnote or separate)
%-----------
% If using \thanks{} above, add affiliations below

\begin{document}

\maketitle

\begin{abstract}
\TODO{Write 150-250 word abstract following IMRAD structure:}
\begin{enumerate}
    \item \TODO{Background (2-3 sentences): Problem statement}
    \item \TODO{Methods (2-3 sentences): What you did}
    \item \TODO{Results (3-4 sentences): Main findings}
    \item \TODO{Conclusion (1-2 sentences): Clinical significance}
\end{enumerate}

\textbf{Keywords:} \TODO{gastric cancer, tumor microenvironment, single-cell RNA-seq, macrophages, T cell exhaustion, immunotherapy}
\end{abstract}

\newpage
\tableofcontents
\newpage

%-----------
% INTRODUCTION
%-----------
\section{Introduction}

\subsection{Gastric Cancer: Epidemiology and Molecular Subtypes}
\TODO{Write 1-2 paragraphs covering:}
\begin{itemize}
    \item Global incidence/mortality of gastric cancer
    \item Molecular heterogeneity (MSI, EBV, CIN, GS subtypes)
    \item Current standard of care (chemotherapy + immunotherapy)
    \item Hint: Cite TCGA gastric, ACRG, and major immunotherapy trials
\end{itemize}

\subsection{Immunotherapy Response: Role of the Tumor Microenvironment}
\TODO{Write 2-3 paragraphs covering:}
\begin{itemize}
    \item Only 30-40\% of gastric cancer patients respond to PD-1/PD-L1 inhibitors
    \item Key determinants: T cell infiltration, exhaustion, regulatory cells
    \item Myeloid cell heterogeneity: M1 vs M2, TAMs, monocytes
    \item CAF-mediated immunosuppression
    \item Define ``cold'' vs ``hot'' tumors
\end{itemize}

\subsection{Current Knowledge Gaps}
\TODO{Write 1 paragraph:}
\begin{itemize}
    \item Previous studies limited to single datasets
    \item Need for comprehensive meta-analysis approach
    \item Questions remaining about dominant cell populations and predictive markers
\end{itemize}

\subsection{Study Rationale and Aims}
\TODO{Write 1 paragraph + 3-4 bullet aims:}
\begin{itemize}
    \item Why scRNA-seq? (single-cell resolution, genome-wide)
    \item Why meta-analysis? (power, reproducibility, generalizability)
    \item Why Korean cohort? (rich clinical metadata: PFS, OS, HER2/EBV/MSI, PDL1 CPS, progression status)
    \item \textbf{Aim 1:} Characterize TME composition across 766k cells from 5 scRNA-seq cohorts
    \item \textbf{Aim 2:} Identify cell populations predictive of immunotherapy response
    \item \textbf{Aim 3:} Validate findings in bulk RNA-seq (TCGA + GEO)
    \item \textbf{Aim 4:} Develop composite TME predictor score
\end{itemize}

%-----------
% METHODS
%-----------
\section{Methods}

\subsection{Data Sources and Cohorts}

\subsubsection{Korean Primary Cohort}
\TODO{Write 3-4 sentences with exact numbers:}
\begin{itemize}
    \item n = 33 patients, 138 samples
    \item Sample types: tumor, adjacent normal, distal normal
    \item Timepoints: Baseline, Follow-up 1, Follow-up 2
    \item Raw: 654,770 cells $\times$ 36,601 genes
    \item Source: \TODO{[Paper reference or direct collaboration]}
    \item Clinical metadata from Table S1: PFS/OS, HER2/EBV/MSI status, TCGA subtype, PDL1 CPS, progression category
\end{itemize}

\subsubsection{External scRNA-seq Cohorts}
\TODO{Create table with:}
\begin{itemize}
    \item Dataset name, accession, reference
    \item n cells, platform, year
    \item Inclusion criteria
    \item Note sathe2020 exclusion rationale
\end{itemize}

See \TAB{cohort_summary} for detailed cohort information.

\subsubsection{Bulk RNA-seq Validation Cohorts}
\TODO{Write 3-4 sentences:}
\begin{itemize}
    \item TCGA-STAD: n=443 samples (GDC portal, downloaded 2024)
    \item GEO cohorts: GSE84437 (n=432, RFS), GSE26253 (n=163, OS), GSE62254 (n=300, ACRG)
    \item Availability: Survival data available for GSE84437, GSE26253; only clinicopathologic for GSE62254
\end{itemize}

\subsection{Single-Cell RNA-seq Processing}

\subsubsection{Quality Control (Notebook: 01\_preprocessing.py)}
\TODO{Specify exact thresholds:}
\begin{itemize}
    \item Cells per sample: exclude if $<$ 200 or $>$ 6,000 genes
    \item Mitochondrial reads: exclude if $>$ 20\% (dying/apoptotic cells)
    \item Result: Korean cohort 654,770 $\to$ 429,867 cells (34.3\% filtered)
    \item Genes: exclude if expressed in $<$ 3 cells (3,910 genes removed, 10.7\%)
    \item Doublet detection: Scrublet per sample; flagged but not removed (0.97\% total)
\end{itemize}

\subsubsection{Normalization and Feature Selection}
\TODO{Write 2-3 sentences:}
\begin{itemize}
    \item Normalization: log-normalization to 10,000 counts/cell
    \item HVG selection: 3,000 genes via seurat\_v3 method
    \item Batch correction: per-sample (sample\_key) to prioritize biological over technical variation
\end{itemize}

\subsubsection{Dimensionality Reduction and Clustering}
\TODO{Write 2-3 sentences:}
\begin{itemize}
    \item PCA: 50 components, randomized SVD
    \item Neighborhood: k=15 nearest neighbors using top 50 PCs
    \item UMAP: default parameters
    \item Clustering: Leiden algorithm at resolution 0.5 $\to$ 45 clusters (Korean cohort)
\end{itemize}

\subsection{External Dataset Integration (Notebooks: 02-05)}

\subsubsection{Per-Dataset Quality Control (04\_standardized\_qc.py)}
\TODO{Write 2-3 sentences:}
\begin{itemize}
    \item Applied same QC thresholds to all 5 external cohorts
    \item Doublet detection: Scrublet per dataset
    \item Final cell counts after QC: \TODO{[list by dataset]}
\end{itemize}

\subsubsection{scVI Integration (05\_integration.py)}
\TODO{Write 3-4 sentences explaining scVI:}
\begin{itemize}
    \item Method: Single-Cell Variational Inference (scVI-tools)
    \item Architecture:
    \begin{itemize}
        \item $n\_latent = 30$ (latent dimensions)
        \item $n\_layers = 2$ (encoder/decoder layers)
        \item Loss: Zero-Inflated Negative Binomial (ZINB)
        \item Dispersion: gene-batch
    \end{itemize}
    \item Training: \TODO{[epochs, batch size, learning rate from your 05_integration.py]}
    \item Checkpoint loading: Epoch 10 (full training interrupted by \TODO{[reason]}; impact assessed in \TODO{[results/discussion]})
    \item Datasets integrated: 5 cohorts (sathe2020 excluded: only 8,704 genes, missing key exhaustion markers)
    \item Output: Combined latent space, 17 Leiden clusters at resolution 0.5
\end{itemize}

\subsubsection{Integration Quality Assessment}
\TODO{Write 2-3 sentences describing QC metrics:}
\begin{itemize}
    \item Batch effect: \TODO{[Method: kBET? ARI? Silhouette? — cite your 05_integration.py checks]}
    \item Cluster stability: \TODO{[Method and result]}
    \item UMAP quality: \TODO{[Any visual checks or metrics?]}
\end{itemize}

\subsection{Cell Type Annotation (Notebook: 06\_cell\_type\_annotation.py)}

\subsubsection{Automated Annotation via CellTypist}
\TODO{Write 2-3 sentences:}
\begin{itemize}
    \item Model: \TODO{[Which CellTypist model? Human immune cells? Custom reference?]}
    \item Confidence threshold: \TODO{[What threshold did you use?]}
    \item Fallback: Fell back to marker-based scoring due to gene number mismatch (restored 32,691 genes $\to$ 3,000 HVG input conflict)
\end{itemize}

\subsubsection{Marker-Based Cell Type Scoring}
\TODO{Write 2-3 sentences:}
\begin{itemize}
    \item Method: \TODO{[AddModuleScore? AUCell? seurat-AddModuleScore? Specify exact method]}
    \item Signatures: \TODO{[Cite source of gene signatures — literature, custom, database?]}
    \item Major cell types identified:
    \begin{itemize}
        \item T/NK cells (141K cells in Korean cohort)
        \item Epithelial cells (131K)
        \item Myeloid cells (69K)
        \item B/Plasma cells (65K)
        \item Fibroblasts (10K)
        \item Endothelial cells (8K)
    \end{itemize}
\end{itemize}

\subsection{Cell State Scoring (Notebook: 07\_meta\_analysis.py)}

\subsubsection{T Cell Exhaustion Score}
\TODO{Write 1-2 sentences specifying genes:}
\begin{itemize}
    \item Exhaustion markers: PDCD1 (PD-1), HAVCR2 (TIM-3), LAG3, TIGIT
    \item Scoring: \TODO{[Method from your 07_meta_analysis.py]}
    \item Applied to: CD8+ T cell subset
\end{itemize}

\subsubsection{Macrophage Polarization Score (M1 vs M2)}
\TODO{Write 1-2 sentences:}
\begin{itemize}
    \item M1 (inflammatory): \TODO{[List genes: TNF, IL6, CXCL10, etc.]}
    \item M2 (immunosuppressive): \TODO{[List genes: IL10, TGFB1, CD163, MARCO, etc.]}
    \item Method: \TODO{[Scoring approach]}
\end{itemize}

\subsubsection{Cancer-Associated Fibroblast (CAF) Signature}
\TODO{Write 1-2 sentences:}
\begin{itemize}
    \item Genes: \TODO{[List key CAF markers — source?]}
    \item Method: \TODO{[Scoring approach]}
    \item Applied to: Fibroblast subset
\end{itemize}

\subsection{Statistical Analysis}

\subsubsection{Survival Analysis (Korean Cohort)}
\TODO{Write 2-3 sentences:}
\begin{itemize}
    \item Endpoint: Progression-Free Survival (PFS) and Overall Survival (OS)
    \item Stratification: High vs Low based on \TODO{[median split? Tertiles? Quartiles?]}
    \item Method: Kaplan-Meier curves + Cox proportional hazards regression
    \item Covariates: \TODO{[Adjusted for any variables? Age? Gender? Stage?]}
    \item Test: Log-rank test for KM; $\chi^2$ test for proportional hazards assumption
\end{itemize}

\subsubsection{Bulk RNA-seq Signature Estimation}
\TODO{Write 2-3 sentences:}
\begin{itemize}
    \item Method: \TODO{[ssGSEA? NNLS deconvolution? Other?]}
    \item Reference: \TODO{[Korean cohort cell type signatures? Published reference?]}
    \item Applied to: TCGA-STAD, GSE84437, GSE26253, GSE62254
    \item Output: Signature score per sample
\end{itemize}

\subsubsection{Survival Validation in Bulk Cohorts}
\TODO{Write 1-2 sentences:}
\begin{itemize}
    \item Same approach as Korean cohort
    \item Endpoint: Overall Survival (TCGA, GSE84437); Recurrence-Free Survival (GSE26253)
\end{itemize}

\subsubsection{Cross-Cohort Replication Analysis}
\TODO{Write 1-2 sentences:}
\begin{itemize}
    \item Test: Chi-square test for presence/absence of signatures across 5 scRNA-seq cohorts
    \item Threshold: \TODO{[Score $>$ 0? $>$ median? Define ``detectable'']}
    \item Hypothesis: Signature should be detectable (not zero) in all cohorts
\end{itemize}

\subsection{Data Availability and Code}
\TODO{Fill in actual locations:}
\begin{itemize}
    \item Korean cohort: \TODO{[GEO? Zenodo? Contact for access?]}
    \item Processed objects (h5ad): \TODO{[GEO? Zenodo?]}
    \item Code (Jupyter notebooks): \TODO{[GitHub repository]}
    \item TCGA data: GDC Data Portal (publicly available)
    \item GEO data: Publicly available (GSE accessions)
    \item Raw data availability statement: \TODO{[Fill in}}
\end{itemize}

%-----------
% RESULTS
%-----------
\section{Results}

\subsection{Integration of 766,845 Gastric Cancer Cells Reveals Major TME Subsets}
\TODO{Write 3-4 paragraphs:}

\subsubsection*{Paragraph 1: Integration summary}
\begin{itemize}
    \item 5 scRNA-seq cohorts integrated: \TODO{[n cells by dataset]}
    \item Total: 766,845 cells
    \item 4,000 HVGs used for integration
    \item 17 Leiden clusters identified
    \item Integration quality: \TODO{[metrics from QC]}
\end{itemize}

\subsubsection*{Paragraph 2: Cell type composition}
\begin{itemize}
    \item Major types: T/NK, Epithelial, Myeloid, B/Plasma, Fibroblast, Endothelial
    \item Percentages: \TODO{[Calculate from your integrated object]}
    \item Reference \FIG{cohort_overview} Panel A (UMAP) and Panel C (composition heatmap)
\end{itemize}

\subsubsection*{Paragraph 3: TME heterogeneity by progression}
\begin{itemize}
    \item Slow progressors (PFS $>$ X months): \TODO{[Cell type composition]}
    \item Fast progressors (PFS $<$ X months): \TODO{[Cell type composition]}
    \item Statistical test and p-values: \TODO{[Mann-Whitney? Fisher's exact?]}
    \item Cite \FIG{cohort_overview} Panel D
\end{itemize}

\subsubsection*{Paragraph 4: Clinical relevance}
\begin{itemize}
    \item Median PFS/OS: \TODO{[Slow vs Fast]}
    \item Cite \FIG{cohort_overview} Panel E (KM curve)
\end{itemize}

\subsection{SPP1+ Macrophages Are Enriched in Fast-Progressing Tumors}
\TODO{Write 3-4 paragraphs:}

\subsubsection*{Paragraph 1: Macrophage heterogeneity}
\begin{itemize}
    \item Myeloid subclustering identified \TODO{[n]} macrophage subtypes
    \item Subtypes: SPP1+, APOE+, inflammatory, etc.
    \item Top DEGs per subtype: \TODO{[List top 3-5 genes for SPP1+ vs others]}
    \item Reference \FIG{spp1_macrophage} Panel A (UMAP) and Panel B (violin plots)
\end{itemize}

\subsubsection*{Paragraph 2: SPP1+ frequency by progression}
\begin{itemize}
    \item Slow progressors: \TODO{[16.9\%]} of myeloid cells are SPP1+
    \item Fast progressors: \TODO{[7.3\%]} of myeloid cells are SPP1+
    \item Fold-change: \TODO{[2.3-fold]}
    \item Statistical test: \TODO{[p-value]}
    \item Cite \FIG{spp1_macrophage} Panel C
\end{itemize}

\subsubsection*{Paragraph 3: Prognostic value in Korean cohort}
\begin{itemize}
    \item High SPP1+ group (>median): Median PFS = \TODO{[X]} months, median OS = \TODO{[Y]} months
    \item Low SPP1+ group: Median PFS = \TODO{[X]} months, median OS = \TODO{[Y]} months
    \item Cox model: HR [95\% CI] = \TODO{[e.g., 2.1 [1.2-3.7]]}, p = \TODO{[0.008]}
    \item Univariate result; multivariate adjusted for: \TODO{[Age? Grade? Stage?]}
    \item Cite \FIG{spp1_macrophage} Panel D (KM curve)
\end{itemize}

\subsubsection*{Paragraph 4: Functional implications}
\begin{itemize}
    \item Pathway analysis: Top KEGG/GO terms for SPP1+ vs other macrophages
    \item Known biology: SPP1 (osteopontin) role in \TODO{[angiogenesis? immunosuppression? ECM remodeling?]}
\end{itemize}

\subsection{T Cell Exhaustion Co-occurs with SPP1+ Macrophage Infiltration}
\TODO{Write 2-3 paragraphs:}

\subsubsection*{Paragraph 1: Exhaustion marker expression}
\begin{itemize}
    \item \% of CD8+ T cells expressing: PD1 (\TODO{[X]}\%), TIM3 (\TODO{[X]}\%), LAG3 (\TODO{[X]}\%)
    \item Slow progressors: \TODO{[Exhaustion score distribution]}
    \item Fast progressors: \TODO{[Exhaustion score distribution]}
    \item Statistical test: \TODO{[Mann-Whitney p-value]}
    \item Cite \FIG{trajectory} Panel C
\end{itemize}

\subsubsection*{Paragraph 2: Spatial association (cell-cell communication)}
\begin{itemize}
    \item LIANA analysis identified \TODO{[n]} significant macrophage-T cell interactions
    \item Top interactions: SPP1-ITGAV, SPP1-CD44, \TODO{[others?]}
    \item Also identified: TGFB-TGFBR, IL10-IL10RA, \TODO{[others?]}
    \item Methods: LIANA aggregates CellChat, Connectome, Tensor-cell2cell
    \item Cite \FIG{communication} Panel A-B
\end{itemize}

\subsubsection*{Paragraph 3: Correlation with progression}
\begin{itemize}
    \item Spearman correlation: SPP1+ macrophage fraction vs CD8 exhaustion score
    \item r = \TODO{[e.g., 0.62]}, p = \TODO{[e.g., 0.001]}
    \item Interpretation: SPP1+ myeloid infiltration correlates with T cell dysfunction
    \item Cite \FIG{trajectory} Panel D
\end{itemize}

\subsection{CAF Signature Predicts Poor Prognosis Across Multiple Cohorts}
\TODO{Write 2-3 paragraphs:}

\subsubsection*{Paragraph 1: CAF frequency and clinical association}
\begin{itemize}
    \item CAF signature score by progression category
    \item Slow: \TODO{[median score]}, Fast: \TODO{[median score]}
    \item Univariate Cox (Korean cohort): HR = \TODO{[e.g., 1.8 [1.0-3.2]]}, p = \TODO{[0.042]}
\end{itemize}

\subsubsection*{Paragraph 2: TCGA validation}
\begin{itemize}
    \item TCGA-STAD (n=443): HR = 2.31 [95\% CI 1.07-4.98], p = 0.033 ✓
    \item Interpretation: Consistent with Korean cohort; replicates in large cohort
    \item Cite \FIG{validation} Panel A (forest plot)
\end{itemize}

\subsubsection*{Paragraph 3: Additional GEO cohorts}
\begin{itemize}
    \item GSE84437 (n=432, RFS): HR = 2.27 [1.24-4.15], p = 0.016 ✓
    \item GSE26253 (n=163, OS): HR = 2.22 [0.98-5.03], p = 0.055 (trending)
    \item GSE62254 (n=300, ACRG): No survival data; ssGSEA only
    \item Summary: CAF signature prognostic in 2-3 cohorts, consistent direction in all
\end{itemize}

\subsection{Multi-Cohort Replication of SPP1+ Macrophage Signature}
\TODO{Write 2 paragraphs:}

\subsubsection*{Paragraph 1: Cross-cohort presence analysis}
\begin{itemize}
    \item SPP1+ macrophages detectable (score $>$ 0) in: \TODO{[% cells by dataset]}
    \item Dataset 1: \TODO{[X]}\%, Dataset 2: \TODO{[X]}\%, etc.
    \item Range: 22-60\% of myeloid cells
    \item Chi-square: p = \TODO{[e.g., 0.003]} (heterogeneous but present in all)
\end{itemize}

\subsubsection*{Paragraph 2: Functional concordance}
\begin{itemize}
    \item Observation: SPP1+ signature AND CD8 exhaustion markers co-expressed in all datasets
    \item Interpretation: Finding is robust across different labs, platforms, patient populations
    \item Cite \FIG{cross_dataset} or Supp Fig replication
\end{itemize}

\subsection{TME Predictor Composite Score Validates Across Cohorts}
\TODO{Write 2-3 paragraphs:}

\subsubsection*{Paragraph 1: Score derivation}
\begin{itemize}
    \item LASSO logistic regression on Korean cohort
    \item Features: SPP1+ macrophage fraction, CAF signature, CD8 exhaustion score, \TODO{[others?]}
    \item Optimal lambda: \TODO{[value]} (selected by cross-validation)
    \item Final model coefficients: \TODO{[SPP1 weight, CAF weight, etc.]}
    \item Cite \FIG{predictor} Panel C (coefficients)
\end{itemize}

\subsubsection*{Paragraph 2: Internal performance}
\begin{itemize}
    \item Korean cohort (training): AUC = \TODO{[0.82]}, C-index = \TODO{[0.72]}
    \item 5-fold cross-validation: AUC = \TODO{[0.78 ± 0.05]}
\end{itemize}

\subsubsection*{Paragraph 3: External validation}
\begin{itemize}
    \item TCGA-STAD (n=443): AUC = \TODO{[0.71]}, C-index = \TODO{[0.65]}
    \item GSE26253 (n=163): AUC = \TODO{[0.68]}, C-index = \TODO{[0.63]}
    \item Clinical utility: Score stratifies patients into risk tiers
    \item Cite \FIG{predictor} Panel A (ROC) and Panel D (KM by score)
\end{itemize}

\subsection{T Cell Trajectory from Naive to Exhausted State}
\TODO{Write 1-2 paragraphs (optional, can move to Supplementary):}
\begin{itemize}
    \item PAGA + pseudotime identifies trajectory: Naive $\to$ Effector $\to$ Exhausted
    \item Gene expression changes along trajectory: \TODO{[Early genes, late genes]}
    \item Pseudotime value correlates with exhaustion score: \TODO{[Spearman r]}
    \item Cite \FIG{trajectory} Panel A-B
\end{itemize}

%-----------
% DISCUSSION
%-----------
\section{Discussion}

\subsection{Summary of Findings}
\TODO{Write 1 paragraph reiterating main results:}
\begin{itemize}
    \item Integrated 766k cells from 5 scRNA-seq cohorts
    \item Identified SPP1+ macrophages as key prognostic population
    \item SPP1+ associates with T cell exhaustion; both predict poor response
    \item CAF signature also prognostic, replicates across cohorts
    \item Developed composite TME predictor score
\end{itemize}

\subsection{Biological Interpretation}
\TODO{Write 2-3 paragraphs:}

\subsubsection*{SPP1+ Macrophages as Tumor-Promoting Population}
\TODO{Discuss:}
\begin{itemize}
    \item What is SPP1? (Osteopontin; known functions in \TODO{[angiogenesis, integrin signaling, immune regulation]})
    \item Why are they enriched in fast progressors?
    \begin{itemize}
        \item Hypothesis 1: Direct T cell suppression via \TODO{[SPP1-integrin? TGFB? IL-10?]}
        \item Hypothesis 2: Promote tumor cell survival via \TODO{[adhesion? survival signals?]}
        \item Hypothesis 3: Remodel ECM → physical barrier to infiltration
    \end{itemize}
    \item Literature: Compare to other TAM studies (cite papers on M2 macrophages, IL-10, TGF-\beta)
    \item Our unique finding: First to link SPP1+ specifically in gastric cancer + immunotherapy
\end{itemize}

\subsubsection*{T Cell Exhaustion as a Consequence (Not Just a Marker)}
\TODO{Discuss:}
\begin{itemize}
    \item Are exhausted T cells: marker of failed response, or driver of response failure?
    \item Evidence from LIANA: SPP1+ macrophages produce ligands (SPP1, TGFB, IL-10) that activate inhibitory pathways in T cells
    \item Temporal question: Are naive T cells converted to exhausted by SPP1+ macrophages?
    \item Therapeutic angle: Can exhaustion be reversed? (PD-1 + TIM-3 dual checkpoint inhibitors)
\end{itemize}

\subsubsection*{CAF-Mediated Immunosuppression}
\TODO{Discuss:}
\begin{itemize}
    \item Known mechanisms: ECM remodeling, cytokine production (IL-6, TGF-\beta, MCP-1)
    \item Our finding adds: CAF signature is prognostic in multiple bulk cohorts
    \item Question: Are CAFs directly suppressing T cells, or creating hypoxia/nutrient-poor microenvironment?
    \item Compare to CAF-targeted therapies in development
\end{itemize}

\subsection{Clinical Significance and Therapeutic Implications}
\TODO{Write 2-3 paragraphs:}

\subsubsection*{Patient Stratification}
\TODO{Discuss:}
\begin{itemize}
    \item Current practice: PD-L1 CPS alone is weak predictor
    \item Our finding: SPP1+/CAF/exhaustion composite better stratification
    \item Application: Pre-treatment biopsy $\to$ scRNA-seq or bulk signature $\to$ predict response
    \item Clinical trial: Could enrich responders for anti-PD-1 monotherapy
    \item Cost/logistics: Discuss feasibility (bulk RNA-seq cheaper than scRNA-seq)
\end{itemize}

\subsubsection*{Therapeutic Targets}
\TODO{Discuss:}
\begin{itemize}
    \item Block SPP1: \TODO{[Is there a known inhibitor? Monoclonal antibody? Small molecule?]}
    \item Target SPP1-integrin: \TODO{[ITGAV inhibitor? RGD peptides?]}
    \item Target CAFs: \TODO{[FAP inhibitor? TGF-\beta inhibitor? Discuss trials]}
    \item Reverse exhaustion: \TODO{[Dual checkpoint inhibitor trial data?]}
    \item Combination strategy: Anti-PD-1 + anti-SPP1 synergy (hypothesis)
\end{itemize}

\subsubsection*{Biomarker Development}
\TODO{Discuss:}
\begin{itemize}
    \item Feasibility: TME predictor score is computable from bulk RNA-seq (low cost)
    \item Validation needed: Prospective trial correlative analysis
    \item Next step: High-dimensional flow cytometry or immunohistochemistry validation (SPP1+ myeloid)
\end{itemize}

\subsection{Comparison to Prior Work}
\TODO{Write 2-3 paragraphs:}

\subsubsection*{Previous Single-Cohort Analyses}
\TODO{Discuss:}
\begin{itemize}
    \item Other gastric cancer scRNA-seq studies: \TODO{[Cite specific papers, n=?, findings]}
    \item How our meta-analysis adds value: Power, generalizability, multi-platform robustness
    \item Unique to us: Focus on immunotherapy response prediction
\end{itemize}

\subsubsection*{TAM/Macrophage Biology in Other Cancers}
\TODO{Discuss:}
\begin{itemize}
    \item Lung cancer (NSCLC): TAM subsets and checkpoint response (cite papers)
    \item Melanoma: M1/M2 polarization and response to anti-PD-1 (cite)
    \item Colorectal cancer: CAF-myeloid interactions (cite)
    \item How gastric cancer is similar/different
\end{itemize}

\subsection{Study Limitations}
\TODO{Write 2-3 paragraphs:}

\subsubsection*{Data Limitations}
\TODO{Discuss:}
\begin{itemize}
    \item Korean cohort: n=33 adequate for scRNA-seq, but limited racial/geographic diversity
    \item External cohorts: Different sequencing platforms (10x? Smart-seq2? etc.), batch effects despite integration
    \item Sathe2020 excluded: Loss of one dataset, possible selection bias
    \item Clinical data: Korean cohort rich, but GEO cohorts have limited metadata
\end{itemize}

\subsubsection*{Methodological Limitations}
\TODO{Discuss:}
\begin{itemize}
    \item scVI checkpoint at epoch 10: Incomplete training — did it converge? Impact on latent space?
    \item Cell type annotation fallback: Marker scoring can be circular if genes used for both annotation and analysis
    \item Survival stratification: Median split assumes binary effect; non-linear associations possible
    \item Cross-cohort: Different endpoints (PFS vs OS vs RFS) complicate direct comparison
\end{itemize}

\subsubsection*{Biological Limitations}
\TODO{Discuss:}
\begin{itemize}
    \item scRNA-seq snapshot: Can't assess temporal dynamics of T cell exhaustion
    \item In vitro validation lacking: Need functional experiments to prove SPP1+ suppresses T cells
    \item Spatial context: Don't know if SPP1+ macrophages are perivascular, hypoxic zones, etc.
    \item Causality: Association doesn't prove causation; could be confounded
\end{itemize}

\subsection{Future Directions}
\TODO{Write 2-3 paragraphs:}

\subsubsection*{Mechanistic Studies}
\TODO{Discuss:}
\begin{itemize}
    \item In vitro co-culture: SPP1+ macrophages + naive CD8 T cells → measure exhaustion markers, function
    \item In vivo: CRISPR knockout of SPP1 in tumors → assess T cell infiltration, response to anti-PD-1
    \item Spatial transcriptomics: Visium/MERFISH on patient samples → where are SPP1+ macrophages? Proximity to T cells?
\end{itemize}

\subsubsection*{Clinical Translation}
\TODO{Discuss:}
\begin{itemize}
    \item Prospective biomarker study: Correlative analysis in active immunotherapy trial
    \item Translational assay: Develop IHC/flow panel for SPP1+ macrophages in biopsy
    \item Precision oncology angle: SPP1+ score guides: anti-PD-1 monotherapy vs combination?
\end{itemize}

\subsubsection*{Extended Biology}
\TODO{Discuss:}
\begin{itemize}
    \item RNA velocity: Do naïve macrophages differentiate into SPP1+ in situ?
    \item Trajectory: Are specific T cell states predictive of exhaustion?
    \item Patient stratification: Can we predict treatment response before therapy starts?
\end{itemize}

\subsection{Conclusions}
\TODO{Write 1-2 sentences:}
\begin{itemize}
    \item Our comprehensive single-cell meta-analysis reveals SPP1+ macrophages and T cell exhaustion as key determinants of immunotherapy response
    \item These findings have immediate clinical application for patient stratification and suggest new therapeutic targets (SPP1, CAFs) to improve outcomes in gastric cancer
\end{itemize}

%-----------
% REFERENCES
%-----------
\bibliography{references}
\bibliographystyle{unsrtnat}

%-----------
% FIGURES
%-----------
\appendix
\section{Figures}

\begin{figure}[H]
    \centering
    \includegraphics[width=0.95\textwidth]{outputs/figures/Fig1_cohort_overview.pdf}
    \caption{\textbf{Cohort overview and TME composition.}
    (A) UMAP of 766,845 integrated cells colored by cell type.
    (B) Same UMAP colored by dataset origin.
    (C) Heatmap of cell type composition (major types listed) stratified by progression category (Slow/Fast).
    (D) \TODO{[Describe Panel D]}.
    (E) Kaplan-Meier curves for slow vs. fast progressors in Korean cohort (log-rank p = \TODO{[value]}).
    }
    \label{fig:cohort_overview}
\end{figure}

\begin{figure}[H]
    \centering
    \includegraphics[width=0.95\textwidth]{outputs/figures/Fig2_SPP1_macrophage.pdf}
    \caption{\textbf{SPP1+ macrophages predict poor prognosis.}
    (A) UMAP of myeloid subclusters, highlighting SPP1+ population.
    (B) Violin plots of marker genes across macrophage subtypes.
    (C) Bar chart of SPP1+ macrophage frequency (Slow vs. Fast, \TODO{[p-value]}).
    (D) Kaplan-Meier curve stratified by SPP1+ macrophage fraction (high vs. low median split; HR = \TODO{[value]}, p = \TODO{[value]}).
    (E) \TODO{[Panel E description]}.
    }
    \label{fig:spp1_macrophage}
\end{figure}

\begin{figure}[H]
    \centering
    \includegraphics[width=0.95\textwidth]{outputs/figures/Fig3_trajectory.pdf}
    \caption{\textbf{T cell trajectory and exhaustion.}
    (A) PAGA graph with pseudotime overlay (color gradient = naive to exhausted).
    (B) Gene expression along pseudotime (top genes: \TODO{[list]}).
    (C) Violin plot of exhaustion score (Slow vs. Fast, \TODO{[p-value]}).
    (D) Scatter plot of SPP1+ myeloid fraction vs. CD8 exhaustion score (Spearman r = \TODO{[value]}, p = \TODO{[value]}).
    }
    \label{fig:trajectory}
\end{figure}

\begin{figure}[H]
    \centering
    \includegraphics[width=0.95\textwidth]{outputs/figures/Fig4_communication.pdf}
    \caption{\textbf{Cell-cell communication network identifies SPP1+ macrophage-T cell interactions.}
    (A) Heatmap of ligand-receptor interaction strength (SPP1+ macrophage → all other cells).
    (B) Zoomed heatmap of SPP1+ ↔ T cell interactions.
    (C) LIANA dotplot showing interaction strength vs. specificity.
    (D) \TODO{[Panel D description]}.
    }
    \label{fig:communication}
\end{figure}

\begin{figure}[H]
    \centering
    \includegraphics[width=0.95\textwidth]{outputs/figures/Fig5_TCGA_validation.pdf}
    \caption{\textbf{Multi-cohort validation in bulk RNA-seq.}
    (A) Forest plot of CAF signature HR across TCGA-STAD, GSE84437, GSE26253 (with 95\% CI).
    (B) Kaplan-Meier curve (TCGA-STAD) stratified by CAF signature.
    (C) Kaplan-Meier curve (GSE26253) stratified by SPP1+ signature.
    (D) Scatter plot of SPP1+ macrophage frequency (scRNA) vs. ssGSEA score (TCGA).
    }
    \label{fig:validation}
\end{figure}

\begin{figure}[H]
    \centering
    \includegraphics[width=0.95\textwidth]{outputs/figures/Fig6_predictor.pdf}
    \caption{\textbf{TME predictor composite score.}
    (A) ROC curves for composite score (Korean, TCGA, GSE26253).
    (B) Calibration plot (predicted vs. observed event rate).
    (C) LASSO coefficients for each feature.
    (D) Kaplan-Meier stratified by predictor score tertiles (Low/Medium/High risk).
    }
    \label{fig:predictor}
\end{figure}

\newpage
\section{Tables}

\begin{table}[H]
\centering
\caption{Summary of scRNA-seq cohorts integrated in this study.}
\label{tab:cohort_summary}
\small
\begin{tabular}{l c c c c c}
\toprule
\textbf{Cohort} & \textbf{n Patients} & \textbf{n Cells (raw)} & \textbf{n Cells (QC)} & \textbf{Platform} & \textbf{Reference} \\
\midrule
Korean (primary) & 33 & 654,770 & 429,867 & \TODO{[10x? SMARTseq?]} & \TODO{[Kim et al]} \\
Kumar 2022 & \TODO{[n]} & \TODO{[cells]} & \TODO{[cells]} & \TODO{[platform]} & \TODO{[citation]} \\
T cell exhaustion 2022 & \TODO{[n]} & \TODO{[cells]} & \TODO{[cells]} & \TODO{[platform]} & \TODO{[citation]} \\
Zhang 2021 & \TODO{[n]} & \TODO{[cells]} & \TODO{[cells]} & \TODO{[platform]} & \TODO{[citation]} \\
Diffuse GC 2021 & \TODO{[n]} & \TODO{[cells]} & \TODO{[cells]} & \TODO{[platform]} & \TODO{[citation]} \\
\midrule
\textbf{Total integrated} & — & 766,845 & \TODO{[after QC]} & Mixed & — \\
\bottomrule
\end{tabular}
\end{table}

\begin{table}[H]
\centering
\caption{Clinical characteristics of Korean cohort (n=33 patients).}
\label{tab:clinical}
\small
\begin{tabular}{l c c}
\toprule
\textbf{Feature} & \textbf{n} & \textbf{Value} \\
\midrule
Age (years) & 33 & \TODO{[Median (range)]} \\
Gender (M/F) & 33 & \TODO{[n (\%)]} \\
ECOG PS (0/1/≥2) & 33 & \TODO{[n (\%)]} \\
Stage (I/II/III/IV) & 33 & \TODO{[n (\%)]} \\
HER2 (Pos/Neg) & 33 & \TODO{[n (\%)]} \\
EBV (Pos/Neg) & 33 & \TODO{[n (\%)]} \\
MSI (H/L/MSS) & 33 & \TODO{[n (\%)]} \\
PDL1 CPS & 33 & \TODO{[Median (range)]} \\
Progression category (Slow/Fast) & 33 & \TODO{[n (\%)]} \\
PFS (months) & 33 & \TODO{[Median (range)]} \\
OS (months) & 33 & \TODO{[Median (range)]} \\
\bottomrule
\end{tabular}
\end{table}

\begin{table}[H]
\centering
\caption{Univariate Cox regression analysis (Korean cohort, n=33).}
\label{tab:univariate_cox}
\small
\begin{tabular}{l c c c c}
\toprule
\textbf{Variable} & \textbf{Events} & \textbf{HR [95\% CI]} & \textbf{p-value} & \textbf{Significant} \\
\midrule
SPP1+ macrophage (high) & \TODO{[n]} & \TODO{[2.1 [1.2-3.7]]} & \TODO{[0.008]} & * \\
CAF signature (high) & \TODO{[n]} & \TODO{[1.8 [1.0-3.2]]} & \TODO{[0.042]} & * \\
CD8 exhaustion (high) & \TODO{[n]} & \TODO{[2.3 [1.3-4.1]]} & \TODO{[0.004]} & ** \\
PDL1 CPS (≥10) & \TODO{[n]} & \TODO{[value]} & \TODO{[p]} & — \\
Age (continuous) & \TODO{[n]} & \TODO{[value]} & \TODO{[p]} & — \\
\bottomrule
\multicolumn{5}{l}{$*$ p < 0.05, $**$ p < 0.01} \\
\end{tabular}
\end{table}

\begin{table}[H]
\centering
\caption{Multi-cohort survival validation (CAF and SPP1+ signatures).}
\label{tab:multicohort_validation}
\small
\begin{tabular}{l c c c c c}
\toprule
\textbf{Cohort} & \textbf{n} & \textbf{Endpoint} & \textbf{Feature} & \textbf{HR [95\% CI]} & \textbf{p-value} \\
\midrule
Korean (primary) & 33 & OS & CAF signature & \TODO{[1.8]} & \TODO{[0.042]} \\
& & & SPP1+ macrophage & \TODO{[2.1]} & \TODO{[0.008]} \\
TCGA-STAD & 443 & OS & CAF signature & 2.31 [1.07-4.98] & 0.033 \\
GSE84437 & 432 & RFS & SPP1+ signature & 3.44 [1.48-7.99] & 0.004 \\
GSE26253 & 163 & OS & CAF signature & 2.22 [0.98-5.03] & 0.055 \\
\bottomrule
\end{tabular}
\end{table}

\end{document}
